% =============================================================================
% EXPERIMENT_PHYSICS_STANDARD_TEMPLATE.tex
%
% CEvNS-pyEFT 新实验物理文档的标准层级。
%
% 固定 subsection：
%   1. 实验目标
%   2. 数据集
%   3. 探测器
%   4. 中微子源
%   5. 振幅计算
%   6. 运动学
%   7. 截面计算
%   8. 事例数计算
%   9. 统计分析
%  10. 计算目标
%
% 原则：
%   - 论文和用户提供的数据决定实验模型。
%   - 项目现有 Response/振幅层决定核物理计算。
%   - 参考程序只用于验证，不作为物理输入。
%   - Slow 先定义物理积分；Fast 只能由 Slow 推导并逐 bin 验证。
%   - 统计分析必须声明输入是原始计数、ON/OFF 数据还是背景扣除 excess。
% =============================================================================

\providecommand{\dd}{\mathrm d}
\providecommand{\eV}{\mathrm{eV}}
\providecommand{\keV}{\mathrm{keV}}
\providecommand{\MeV}{\mathrm{MeV}}
\providecommand{\eVee}{\mathrm{eV_{ee}}}
\providecommand{\keVee}{\mathrm{keV_{ee}}}
\providecommand{\keVnr}{\mathrm{keV_{nr}}}
\providecommand{\cm}{\mathrm{cm}}
\providecommand{\kgday}{\mathrm{kg\,day}}
\providecommand{\code}[1]{\texttt{#1}}

\section{<实验名称>}
\label{sec:<experiment-label>}

\subsection{实验目标}

说明实验研究的入射粒子、靶核、散射过程和最终观测量。
例如
\begin{equation}
  \nu_\alpha + A
  \longrightarrow
  \nu_\beta + A.
\end{equation}
定义真实反冲变量 $T$ 和实验重建变量 $E_{\rm reco}$。

给出本次接入目标：
\begin{enumerate}
  \item 标准模型或指定 EFT 情景的信号预测；
  \item 重建能谱和分箱事例数；
  \item Slow 计算；
  \item Fast 计算；
  \item 指定的统计分析。
\end{enumerate}

\subsection{数据集}

给出唯一的数据集身份：
\begin{equation}
  \code{dataset\_id}
  =
  \code{<unique-dataset-id>}.
\end{equation}

说明：
\begin{enumerate}
  \item 实验地点和运行阶段；
  \item 数据采集时间；
  \item reactor-on、reactor-off 或其他运行模式；
  \item 对应论文、补充材料和公开数据；
  \item 本文件是否描述一个数据集，还是多个必须拆分的数据集。
\end{enumerate}

不同运行阶段、不同探测器配置和不同统计数据不得混入同一数据集。

\subsection{探测器}

\subsubsection{探测器集合、质量和曝光}

给出探测器集合
\begin{equation}
  d\in\{d_1,d_2,\ldots\}.
\end{equation}

对每台探测器给出 active mass、live time 和曝光：
\begin{equation}
  \mathcal E_d
  =
  M_d\,t_d.
\end{equation}

如果直接使用实验给出的 $\mathcal E_d$，
不得再次乘 $M_d$ 或 $t_d$。

\subsubsection{探测器阈值与分析窗口}

对每台探测器给出
\begin{equation}
  E_{{\rm th},d},
  \qquad
  E_{{\rm min},d},
  \qquad
  E_{{\rm max},d},
\end{equation}
以及精确 bin edges
\begin{equation}
  E_0,E_1,\ldots,E_{N_{\rm bin}}.
\end{equation}

\subsubsection{靶材}

给出材料、同位素丰度和每千克核数。
对同位素 $a$，
\begin{equation}
  n_{a/{\rm kg}}
  =
  \frac{N_A\,10^3}{M_{\rm mol}}\eta_a.
\end{equation}

同时给出项目 Response 数据库中对应的 target/response key。
同位素丰度与核响应是两个独立输入。

\subsection{中微子源}

\subsubsection{通量模型}

事件率最终需要探测器位置的微分通量
\begin{equation}
  \frac{\dd\Phi_s}{\dd E_\nu}.
\end{equation}

本节只写当前数据集实际采用的通量模型，不罗列无关的通用模式。
必须给出：
\begin{enumerate}
  \item 源分量 $s$；
  \item 能量范围；
  \item 通量单位；
  \item 谱形公式或数值表；
  \item 归一化；
  \item 插值、拼接和外推规则；
  \item 通量不确定性。
\end{enumerate}

若由每次裂变谱构造探测器通量，
则明确写出
\begin{equation}
  \frac{\dd\Phi}{\dd E_\nu}
  =
  \frac{P_{\rm th}}
  {4\pi L^2\sum_r f_re_r}
  \sum_r f_rS_r(E_\nu).
\end{equation}
若直接使用探测器位置通量，
则不得再次乘热功率或几何因子。

\subsection{振幅计算}

实验模块默认调用 CEvNS-pyEFT 当前的振幅与核响应层。
对同位素 $a$，
\begin{equation}
  \mathcal M_a
  =
  \mathcal M_a
  \left(
    E_\nu,T;
    \boldsymbol C,
    \mathcal R_a
  \right),
\end{equation}
其中 $\boldsymbol C$ 是标准模型或 EFT 参数，
$\mathcal R_a$ 是项目 Response 层给出的核响应。

本节只需说明：
\begin{enumerate}
  \item 使用哪个理论参数基准点；
  \item 使用哪些 Response keys；
  \item 新模型是否仍被现有振幅接口覆盖；
  \item 若不覆盖，理论层需要修改什么。
\end{enumerate}

标准模型弱荷、form factor 或核响应的具体实现，
应合并在本节的一个说明中，不需要拆成多个实验小节。

\subsection{运动学}

对弹性核散射，
\begin{equation}
  E_{\nu,\min}(T)
  =
  \frac{
    T+\sqrt{T^2+2m_AT}
  }{2},
\end{equation}
\begin{equation}
  T_{\max}(E_\nu)
  =
  \frac{2E_\nu^2}{m_A+2E_\nu}.
\end{equation}

列出当前实验使用的真实变量、重建变量、积分上下限和单位约定。
任何近似都必须显式命名，不能静默替换精确运动学。

\subsection{截面计算}

微分截面由项目振幅层统一计算：
\begin{equation}
  \frac{\dd\sigma_a}{\dd T}
  =
  \mathcal F_\sigma
  \left[
    \overline{
      \left|
        \mathcal M_a
      \right|^2
    }
  \right].
\end{equation}

本节给出：
\begin{enumerate}
  \item 截面输入和输出单位；
  \item 同位素求和方式；
  \item source flavor；
  \item 理论参数入口；
  \item 截面与 Response 的映射；
  \item 运动学区间外返回零的规则。
\end{enumerate}

实验文档不另写一套与项目理论层重复的 CEvNS 截面。

\subsection{事例数计算}

\subsubsection{淬灭、分辨率与效率}

给出真实核反冲到平均电子等效能量的映射
\begin{equation}
  \overline E_{\rm ee}(T)
  =
  1000\,Q(T)\,T,
\end{equation}
其中 $T$ 以 $\keVnr$ 输入，
$\overline E_{\rm ee}$ 以 $\eVee$ 输出。

给出重建响应
\begin{equation}
  \mathcal R_d(E_{\rm reco}\mid T)
\end{equation}
和效率
\begin{equation}
  \epsilon_d(E_{\rm reco}).
\end{equation}

必须说明效率作用于真实变量还是重建变量，
并说明公开数据是否已经过效率校正。

\subsubsection{Slow 事例数}

真实反冲谱为
\begin{equation}
  \frac{\dd N_{das}^{\rm true}}{\dd T}
  =
  n_{a/{\rm kg}}\,
  \mathcal E_d
  \int
  \dd E_\nu\,
  \frac{\dd\Phi_s}{\dd E_\nu}
  \frac{\dd\sigma_a}{\dd T}.
\end{equation}

重建谱为
\begin{equation}
  \frac{\dd N_{das}}
  {\dd E_{\rm reco}}
  =
  \epsilon_d(E_{\rm reco})
  \int
  \dd T\,
  \mathcal R_d(E_{\rm reco}\mid T)
  \frac{\dd N_{das}^{\rm true}}{\dd T}.
\end{equation}

第 $i$ 个 bin 的事例数为
\begin{equation}
  N_{dasi}^{\rm slow}
  =
  \int_{E_i}^{E_{i+1}}
  \dd E_{\rm reco}\,
  \frac{\dd N_{das}}
  {\dd E_{\rm reco}}.
\end{equation}

总预测为
\begin{equation}
  N_i^{\rm slow}
  =
  \sum_d\sum_a\sum_s
  N_{dasi}^{\rm slow}.
\end{equation}

\subsubsection{Fast 事例数}

先定义 detector-bin kernel
\begin{equation}
  W_{di}(T)
  =
  \int_{E_i}^{E_{i+1}}
  \dd E_{\rm reco}\,
  \epsilon_d(E_{\rm reco})
  \mathcal R_d(E_{\rm reco}\mid T).
\end{equation}

则
\begin{equation}
  N_i
  =
  \sum_{d,a,s}
  n_{a/{\rm kg}}\,
  \mathcal E_d
  \int\dd T\,
  W_{di}(T)
  \int\dd E_\nu\,
  \frac{\dd\Phi_s}{\dd E_\nu}
  \frac{\dd\sigma_a}{\dd T}.
\end{equation}

若振幅对 Wilson 系数线性，
还可进一步构造
\begin{equation}
  N_i^{\rm fast}(\boldsymbol C)
  =
  \boldsymbol C^\dagger
  K_i
  \boldsymbol C.
\end{equation}

Fast 必须逐 bin 与 Slow 比较，不允许仅比较总事例数。

\subsection{统计分析}

必须先声明统计输入模式，只能选择一个：
\begin{enumerate}
  \item 原始观测计数与背景；
  \item reactor-on/reactor-off；
  \item 已经背景扣除的 excess counts。
\end{enumerate}

\subsubsection{背景扣除 excess 的简化卡方}

若输入为
\begin{equation}
  D_i
  =
  \text{background-subtracted excess},
\end{equation}
则可使用
\begin{equation}
  \chi^2
  =
  \sum_i
  \frac{
    \left[
      D_i-(1+\alpha)N_i^{\rm th}
    \right]^2
  }{
    \sigma_i^2
  }
  +
  \left(
    \frac{\alpha}{\sigma_\alpha}
  \right)^2.
\end{equation}

此时不得再把背景 $B_i$ 加入理论预测。

\subsubsection{观测计数与背景的简化卡方}

若输入为原始观测计数 $O_i$ 和背景预测 $B_i$，
则可使用
\begin{equation}
  \chi^2
  =
  \sum_i
  \frac{
    \left[
      O_i
      -(1+\alpha)N_i^{\rm th}
      -(1+\beta)B_i
    \right]^2
  }{
    \sigma_i^2
  }
  +
  \left(
    \frac{\alpha}{\sigma_\alpha}
  \right)^2
  +
  \left(
    \frac{\beta}{\sigma_\beta}
  \right)^2.
\end{equation}

Mathematica 参考程序采用哪一种形式，
本文件就明确选择哪一种；
参考程序只决定本次统计复现的公式，
不决定 Flux、Response 或探测器物理。

\subsection{计算目标}

必须给出可量化验收目标：
\begin{enumerate}
  \item Flux 归一化；
  \item 各探测器和同位素贡献；
  \item 标准模型总事例数；
  \item 逐 bin Slow 结果；
  \item 逐 bin Fast/Slow 相对误差；
  \item 指定参数点的 $\chi^2$；
  \item 论文或用户提供的 benchmark；
  \item 当前可以声明的复现等级。
\end{enumerate}

不得通过整体缩放强制命中 benchmark。
